\FloatBarrier
\newpage
\section{函数研究的方法}
\label{sec:calculus-function-comparison}

\begingroup
\setlength{\parskip}{3.5pt}

前面已经由导数建立单调性、极值和凸性的判据。实际问题还需要确定研究哪个函数：题设可能给出两个数之间的关系，也可能只说明某个点满足一个方程；所求的量未必能直接算出。本节从这些关系出发，选择研究对象，利用已有性质完成比较。

\subsection{选择研究函数}

若\(H\)在区间上严格递增，则\(H(u)<H(v)\)与\(u<v\)等价，\(H(u)=H(v)\)则给出\(u=v\)。因此，有些等式或不等式可以先整理成同一个函数在两处的取值，再研究这个函数的单调性。

例如，\(u,v>0\)时，关系\(u+\ln u=v+\ln v\)可以写成\(H(u)=H(v)\)，其中\(H(x)=x+\ln x\)。由\(H'(x)=1+1/x>0\)，立即得到\(u=v\)。若各项分别比较，反而容易丢失整个表达式的关系。

比较幂也可以转为比较函数值。对\(u,v>0\)，取对数后，\(u^v<v^u\)等价于\(\ln u/u<\ln v/v\)。因此应研究\(H(x)=\ln x/x\)，其导数为\((1-\ln x)/x^2\)。这个函数在\((0,e)\)上递增、在\((e,+\infty)\)上递减；比较两个幂，便转为确定\(u,v\)所在的区间及相应函数值。

有时题设直接给出的是导数的组合。反过来观察乘积法则、商法则与链式法则，有
\begin{equation*}
 \begin{aligned}
 (e^{-ax}f(x))'&=e^{-ax}(f'(x)-af(x)),\\
 \left(\frac{f(x)}x\right)'&=\frac{xf'(x)-f(x)}{x^2},\\
 (\ln|f(x)|)'&=\frac{f'(x)}{f(x)}.
 \end{aligned}
\end{equation*}
这些函数可以沿求导公式反向寻找。若取\(H=wf\)，则\(H'=wf'+w'f\)；要使它包含\(f'-af\)，便希望\(w'=-aw\)，而\(w=e^{-ax}\)恰好满足。类似地，商的导数自然产生\(xf'-f\)，所以研究\(f/x\)。各式均在表达式有意义的区间上使用。

\begin{example}
设\(f(0)=1\)，且\(x\ge0\)时\(f'(x)\ge f(x)\)。证明\(f(x)\ge e^x\)。再设\(x>0\)时\(xf'(x)-f(x)>0\)，比较\(f(2x)\)与\(2f(x)\)。分别说明为什么选择乘积或商。
\end{example}

若新函数仍不容易判号，可以继续求导，再由已知点逐级回推。例如，若\(H(a)=H'(a)=0\)，且\(H''\ge0\)于\([a,b]\)，则先有\(H'\ge0\)，继而有\(H\ge0\)。初值确定了每一步比较的起点。

\begin{example}
比较\(e^\pi\)与\(\pi^e\)的大小。再设\(e<a<b\)，比较\(a^b\)与\(b^a\)。
\end{example}

\newpage
\begin{exercise}
\mbox{}\par\nobreak

\begin{samepage}
（1）设\(u,v\in\mathbb R\)，且\(u^3+u=-(v^3+v)\)，证明\(u+v=0\)。
\par
\end{samepage}

\begin{samepage}
（2）设\(f\)在\([0,1]\)上连续可导，\(f(0)=2\)，且\(f'+2f\le0\)。估计\(f(x)\)，并说明比较方向怎样由所选函数的导数确定。
\par
\end{samepage}

（3）设\(f>0\)，且\(f'/f\le1\)于\([a,b]\)。证明\(f(b)/f(a)\le e^{b-a}\)，比较研究\(\ln f\)与\(e^{-x}f\)两条路线。

（4）设\(H(0)=H'(0)=H''(0)=0\)，且\(H'''\ge0\)于\([0,1]\)。证明\(H\ge0\)。再将这一推理用于\(e^x-1-x-x^2/2\)，指出每次回推所需的零点信息。
\end{exercise}

\subsection{利用方程中的关系}

方程的根即使难以显式求出，所满足的等式仍能参与运算。把这个等式代入待研究的表达式，往往可以减少未知量，求出目标值或确定它的范围。

考察\(f_a(x)=e^x+x^2-ax\)。其导函数\(e^x+2x-a\)严格递增，两端分别趋于负无穷和正无穷，所以有唯一零点\(c\)，且\(c\)为原函数的最小值点。虽然不便显式写出\(c\)，仍有
\begin{equation*}
 a=e^c+2c,\qquad
 f_a(c)=e^c+c^2-c(e^c+2c)=(1-c)e^c-c^2.
\end{equation*}
原来的函数值含有\(a,c\)，利用驻点方程后只剩\(c\)。若还需限定\(a\)的范围，则应由\(a=e^c+2c\)及其单调性，把\(c\)的范围对应回去。这种设出而不显式求出的零点，常称为隐零点。

对于两个未知点，方程还可以同时约束它们的和与积。若\(c_1,c_2\)是二次方程的根，研究对称式时可以先使用韦达关系；研究单个函数值时，也可以用方程替换其中的二次项。第四章在曲线联立中保留根与系数关系的做法，在这里同样适用。

\begin{example}
设\(f_a(x)=e^x+x^2-ax\)，最小值点为\(c\)。确定\(c>0\)、\(c=0\)、\(c<0\)分别对应的\(a\)的范围，再证明最小值不超过\(1\)，并确定等号条件。比较直接使用\(f_a(0)\)与研究化简后的\(f_a(c)\)两种方法。
\end{example}

保留方程以后，还要注意所设点的范围。导函数的一个零点是否为极值点，取决于两侧符号；所得参数表达式是否可以反过来确定这个点，则取决于相应函数的单调性。方程、位置与函数变化共同组成了推理的依据。

\(G_a(x)=\ln x-ax^2+x\)、\(a>0\)的唯一驻点\(c\)满足\(2ac^2=c+1\)，最大值为\(G_a(c)=\ln c+(c-1)/2\)。后者随\(c\)递增，而\(a=(c+1)/(2c^2)\)随\(c\)递减，故最大值随\(a\)递减。

若要求\(G_a(x)\le M\)对所有\(x>0\)成立，就等价于\(\ln c+(c-1)/2\le M\)。先确定允许的\(c\)，再由驻点关系还原参数范围，恒成立问题便落到一个最值点上。

\begin{exercise}
\mbox{}\par\nobreak

（1）设\(c\)为\(e^x+x-2=0\)的根，证明\(0<c<1\)。利用\(e^c=2-c\)，将\(e^{2c}+2ce^c+c^2\)化为确定的数。

（2）设\(f(x)=x(\ln x-ax)\)有两个极值点\(c_1<c_2\)。由驻点方程化简\(f(c_i)\)，并研究如何利用\(c_i\)的位置估计这两个极值。

（3）设\(a>0\)，函数\(f(x)=-\ln x-ax^2+x\)有两个极值点\(c_1,c_2\)。先确定\(a\)的范围，再用韦达关系将\(f(c_1)+f(c_2)\)表示为\(a\)的函数。

（4）对正文中的\(f_a\)，若\(a\)只在\([1,e+2]\)内变化，确定最小值点的范围。说明先求\(c\)的范围，再研究\((1-c)e^c-c^2\)，怎样避免同时保留两个独立变量。
\end{exercise}

\subsection{减少变量}

表达式中出现两个字母，未必就有两个独立变化的量。若它们满足一个方程，可以先利用约束；若没有这样的约束，则可以固定其中一个，将另一个作为自变量研究。变量的选择需要服从题设中的关系。

例如，设\(0<u<v\)，且\(u-\ln u=v-\ln v\)，则
\begin{equation}\label{eq:calculus-equal-level-log}
 v-u=\ln\frac vu.
\end{equation}
若选择比值\(t=v/u>1\)，由\(u(t-1)=\ln t\)，得到
\begin{equation*}
 u=\frac{\ln t}{t-1},\qquad v=\frac{t\ln t}{t-1}.
\end{equation*}
若选择差\(d=v-u>0\)，则\(v=e^du\)，从而
\begin{equation*}
 u=\frac d{e^d-1},\qquad v=\frac{de^d}{e^d-1}.
\end{equation*}
两种选择都利用了\(v-u=\ln(v/u)\)：比值确定了差，差也确定了比值。两种表示满足\(t=e^d\)；将所求的和、积等表达式代入后，都可以转为单变量问题。

参数有时也可以通过换元从等值关系中消去。设\(k>0\)、\(F_k(x)=x-k\ln x\)，令\(x=kt\)，则
\begin{equation*}
 F_k(kt)=k\bigl(t-\ln t-\ln k\bigr).
\end{equation*}
若\(F_k(u)=F_k(v)\)，令\(u=kr\)、\(v=ks\)，即得\(r-\ln r=s-\ln s\)。参数只引起平移、伸缩；再由\(u+v=k(r+s)\)、\(uv=k^2rs\)还原两点的和、积。

若只给出\(0<u<v\)，则不能擅自加入式\eqref{eq:calculus-equal-level-log}这样的等式。此时可以固定\(u\)，研究\(v\)怎样变化。例如，令\(H(v)=\ln(v/u)-(v-u)/u\)，有\(H(u)=0\)、\(H'(v)=1/v-1/u<0\)，故\(\ln(v/u)<(v-u)/u\)。这里仍需保证所得结论对每个允许的\(u\)成立。

\begin{example}
设\(u<v\)且\(e^u-u=e^v-v\)。令\(d=v-u\)，把\(e^u,e^v\)表示为\(d\)的函数，再研究\(u+v\)。比较差值与比值在本题中的便利程度。
\end{example}

\begin{exercise}
\mbox{}\par\nobreak

（1）对满足式\eqref{eq:calculus-equal-level-log}的\(u<v\)，分别用\(t\)、\(d\)表示\(u+v\)与\(uv\)，核对其对应关系。

（2）对\(0<u<v\)，分别固定\(u\)或令\(t=v/u\)，证明\((v-u)/v<\ln(v/u)<(v-u)/u\)。

（3）设\(0<u<v\)、\(ue^{-u}=ve^{-v}\)。取对数后说明它与式\eqref{eq:calculus-equal-level-log}的关系。

（4）设\(u,v>0\)、\(u+v=s\)。用消元、对称性或对数凹性求\(\ln u+\ln v\)的最大值。
\end{exercise}

\subsection{比较零点的位置}

设函数先严格递减、后严格递增，在\(c\)处取得最小值，而\(u<c<v\)满足\(f(u)=f(v)\)。要比较\(u+v\)与\(2c\)，比较\(u\)与反射点\(2c-v\)。若反射点位于左侧单调区间，原来的数值比较就转成\(f(u)\)与\(f(2c-v)\)的比较；再用\(f(u)=f(v)\)，两者可统一为一个变量。

以\(f(x)=x-\ln x\)为例，最小值点为\(1\)。对\(0<t<1\)，考察
\begin{equation*}
 R(t)=f(1+t)-f(1-t)
 =2t-\ln\frac{1+t}{1-t}.
\end{equation*}
有\(R(0)=0\)、\(R'(t)=-2t^2/(1-t^2)<0\)，故右侧等距离位置的函数值更小。若\(1<v<2\)，则\(f(u)=f(v)<f(2-v)\)，由左侧严格递减，得到\(u>2-v\)；若\(v\ge2\)，则\(u+v>2\)直接成立。因此两根的中点在\(1\)的右侧。

另一条路线是直接利用两根满足的方程。为此先建立一个关于正数的平均值关系。

\begingroup
\setlength{\parskip}{3pt}
\setlength{\abovedisplayskip}{8pt}
\setlength{\belowdisplayskip}{8pt}
\begin{proposition}\label{prop:calculus-logarithmic-mean}
若\(0<a<b\)，其对数平均值满足
\begin{equation}\label{eq:calculus-logarithmic-mean}
 \sqrt{ab}<\frac{b-a}{\ln b-\ln a}<\frac{a+b}{2}.
\end{equation}
\end{proposition}

\begin{samepage}
\begin{proof}
令\(q=b/a>1\)。对\(t>1\)，函数\(t-t^{-1}-2\ln t\)在\(1\)处为零，导数为\((t-1)^2/t^2>0\)。取\(t=\sqrt q\)，得到\(\ln q<(q-1)/\sqrt q\)，即左侧不等式。

函数\(\ln q-2(q-1)/(q+1)\)在\(1\)处为零，导数为\((q-1)^2/[q(q+1)^2]>0\)，得到右侧不等式。
\end{proof}
\end{samepage}
\endgroup
\newpage
若\(u-\ln u=v-\ln v\)，式\eqref{eq:calculus-equal-level-log}使对数平均值恰为\(1\)。于是同时得到\(uv<1\)、\(u+v>2\)。取\(a=e^\alpha,b=e^\beta\)，还可将式\eqref{eq:calculus-logarithmic-mean}写成
\begin{equation*}
 e^{(\alpha+\beta)/2}
 <\frac{e^\beta-e^\alpha}{\beta-\alpha}
 <\frac{e^\alpha+e^\beta}{2}\qquad(\alpha<\beta).
\end{equation*}
这些关系把两点函数值之差、两点间的平均变化与几种平均值联系起来。后面建立积分以后，还可以从凸函数的积分估计重新理解它们。

\begin{example}
设\(f(x)=xe^{-x}\)，\(u<v\)且\(f(u)=f(v)>0\)。先用反射比较证明\(u+v>2\)，再分别用比值换元与式\eqref{eq:calculus-logarithmic-mean}证明\(u+v>2\)、\(uv<1\)。比较三条路线所用的数量关系。
\end{example}

非标准的根关系也可以先改变研究对象。例如，要比较\(u+v\)与一个依赖\(u\)的表达式，先把结论整理为\(v>\varphi(u)\)，再研究\(f(\varphi(u))-f(u)\)。构造后必须确认\(v,\varphi(u)\)位于可比较的同一单调区间。目标式决定比较点，原方程则负责把另一个根的函数值消去。

根和与根积对应的比较点也不同。要证明\(u+v>2c\)，可比较\(v\)与\(2c-u\)；要证明\(uv<K\)且\(u>0\)，则可比较\(v\)与\(K/u\)。把目标改写为两点的位置关系，便为辅助函数的选择提供了依据。

\begin{exercise}
\mbox{}\par\nobreak

（1）若\(u<v\)是\(x-\ln x=a\)的两个根，证明\(a>1\)、\(0<u<1<v\)，并由\(u+v=2a+\ln(uv)\)证明\(u+v<2a\)。

（2）在上一题条件下，证明\(u+v>a+1\)。可将目标改写为\(v>e^{1-u}\)，比较\(f(e^{1-u})\)与\(f(u)\)，其中\(f(x)=x-\ln x\)。

（3）对\(0<a<b\)，用对数平均不等式估计\(\ln b-\ln a\)。比较所得两侧界与中值定理给出的\((b-a)/b\)、\((b-a)/a\)。

（4）若\(f\)在\(c\)两侧分别严格递增、严格递减，且\(f(c+t)>f(c-t)\)，研究同一高度的两点\(u<c<v\)的中点位于\(c\)的哪一侧。写出所用比较点处于定义域的条件。

（5）设\(k>0\)，\(u<v\)满足\(u-k\ln u=v-k\ln v\)。证明\(u+v>2k\)、\(uv<k^2\)，并说明这两个界为什么随\(k\)分别作一次、二次伸缩。

（6）对\(0<t<1\)，证明\(2t<\ln\frac{1+t}{1-t}<\frac{2t}{\sqrt{1-t^2}}\)，并说明它与反射差\(R(t)\)的关系。

（7）（2021年新高考I卷第22题第（2）问）设\(a,b\)是两个不相等的正数，且\(b\ln a-a\ln b=a-b\)。证明\(2<1/a+1/b<e\)。
\end{exercise}

\endgroup
\FloatBarrier
\newpage
\subsection{零点随参数的变化*}

固定参数时，零点是一个数；让参数变化，同一单调支上的零点也可以形成一个函数。先确定这个函数存在，再研究它的变化，能够把一组方程的根联系起来。

若方程可以写成\(a=\varphi(x)\)，且\(\varphi\)在所选区间上连续严格单调，就可以用反函数\(x=x(a)\)描述这一支的根。在\(\varphi'(x)\ne0\)处，由反函数求导公式，
\begin{equation*}
 x'(a)=\frac1{\varphi'(x(a))}.
\end{equation*}
这既说明根怎样随参数移动，也给出移动速度。若研究两个根的和、积，分别确定两支后，再对组合求导。

以\(x-ae^x=0\)为例，写成\(a=xe^{-x}\)。在\(0<a<1/e\)时，两支分别给出\(0<u(a)<1<v(a)\)。对\(x-ae^x=0\)沿每一支求导，利用\(ae^x=x\)，得到
\begin{equation*}
 x'(a)=\frac{e^x}{1-ae^x}=\frac{x}{a(1-x)}.
\end{equation*}
因此左根随\(a\)增大而增大，右根则减小。对于它们的和，有
\begin{equation*}
 (u+v)'=\frac{u+v-2uv}{a(1-u)(1-v)}.
\end{equation*}
分母为负。两根还满足\(u-\ln u=v-\ln v\)，故由前面的平均值关系，\(u+v>2\)、\(uv<1\)，分子为正。于是两根之和随\(a\)增大而减小。

单独观察两根一增一减，并不能决定它们的和怎样变化。零点方程、分支导数和根之间的大小关系在这里分别提供了不同的信息。

\begin{example}
对上述两支根\(u(a),v(a)\)，研究当\(a\to0^+\)及\(a\to(1/e)^-\)时它们各自的极限。再研究\(v(a)/u(a)\)的单调性，比较直接求商的导数与先求对数导数的方法。
\end{example}

\begin{exercise}
\mbox{}\par\nobreak

（1）设\(u(a)<v(a)\)为\(x-\ln x=a\)的两个根。求两支的导数，判断根的间距\(v-u\)怎样随\(a\)变化。

（2）证明上题中的\(u+v\)随\(a>1\)增大而增大，并研究\(a\to1^+\)时的极限。与正文中的参数\(a\)比较：两个问题中的参数是否表示同一个量？

（3）对\(a=e^c+2c\)确定的\(c=c(a)\)，证明\(c'(a)=1/(e^c+2)\)。再对最小值\(m(a)=e^{c(a)}+c(a)^2-ac(a)\)求导，观察驻点方程怎样使含\(c'\)的项消去。
\end{exercise}

\subsection{比较函数的构造*}

前面的切线和Taylor多项式提供了简单的局部模型。若希望比较式在较大范围内仍容易处理，也可以选择多项式之比。选择时先匹配局部变化，再单独证明所需区间内的误差符号。

以\(\ln(1+h)\)在零点附近的二阶展开为例，尝试
\begin{equation*}
 Q(h)=\frac{Ah}{1+Bh}=Ah-ABh^2+o(h^2).
\end{equation*}
匹配\(h-h^2/2\)，得到\(A=1\)、\(B=1/2\)，即\(Q(h)=2h/(2+h)\)。这种通过匹配Taylor系数构造的有理逼近，称为Padé逼近。这里仅研究这个低阶例子。

局部匹配并不自行保证全区间的大小关系。令\(E(h)=\ln(1+h)-2h/(2+h)\)，则
\begin{equation*}
 E(0)=0,\qquad E'(h)=\frac{h^2}{(1+h)(2+h)^2}.
\end{equation*}
由此知道，在\(h>0\)时\(E(h)>0\)，在\(-1<h<0\)时\(E(h)<0\)。换成\(x=1+h\)，得到
\begin{equation*}
 \ln x\begin{cases}
 >2(x-1)/(x+1),&x>1,\\
 <2(x-1)/(x+1),&0<x<1.
 \end{cases}
\end{equation*}
现在这个比较已由区间上的导数符号证明，可以用于该区间中的零点定位。

\paragraph{由比较曲线约束零点} \mbox{}

设\(u<v\)是\(f(x)=x-\ln x=a\)、\(a>1\)的两个根。取比较函数
\begin{equation*}
 g(x)=x-\frac{2(x-1)}{x+1}
     =\frac{x^2-x+2}{x+1}.
\end{equation*}
在\((0,1)\)上\(g<f\)，在\((1,+\infty)\)上\(g>f\)，且\(g\)在这两段分别严格递减、严格递增。方程\(g(x)=a\)转为
\begin{equation*}
 x^2-(a+1)x+2-a=0.
\end{equation*}
它有两个根\(r_-<1<r_+\)。右支比较给出\(v>r_+\)；若\(r_->0\)，左支比较给出\(u>r_-\)，若\(r_-\le0\)，这个不等式也直接成立。因此
\begin{equation*}
 u+v>r_-+r_+=a+1.
\end{equation*}
超越方程不易显式求根，而比较方程是二次方程，韦达关系恰好给出所需的根和。这里先证明曲线的上下关系，再在各单调支上严格比较根的位置。

同一个结论也可以按目标直接构造\(f(\varphi(u))-f(u)\)。两条路线分别利用了目标表达式与简单比较曲线；哪一条更方便，要看能否找到易于判号的差，以及比较方程是否保留所需的和、积信息。

\begin{example}
设\(u<v\)为\(x-\ln x=a\)的两个根。取\(q(x)=(x+x^{-1})/2\)，证明它在两根所在的两段上分别位于\(f\)的上、下方，再由\(q(x)=a\)的根证明\(u+v<2a\)。与利用\(uv<1\)的证明比较。
\end{example}

\begin{exercise}
\mbox{}\par\nobreak

（1）设\(Q(x)=(1+Ax)/(1+Bx)\)。使它在零点匹配\(e^x\)直到二阶变化，求\(A,B\)。再在\((-2,2)\)内研究\(e^x\)与\(Q(x)\)的大小关系，可以先比较\(x\)与\(\ln Q(x)\)。

（2）在\(x=1\)附近比较\(\ln x\)与\(2(x-1)/(x+1)\)的Taylor多项式。两者从哪一阶开始不同？这能否直接代替正文中全区间的误差符号证明？

（3）设\(f,g\)在\(x>c\)上均严格递增，且\(g(x)>f(x)\)。若\(r,s>c\)且\(f(r)=g(s)\)，判断\(r,s\)的大小。再写出严格递减时的结论，用以核对正文两支比较的方向。

（4）在\(x=0\)处，用\((1+Ax)/(1+Bx)\)匹配\(\sqrt{1+x}\)直到二阶变化。研究所得有理函数与\(\sqrt{1+x}\)在\(x\ge0\)时的大小关系，并估计\(\sqrt{1.1}\)。

（5）（2018年全国I卷理科第21题第（2）问）设\(f(x)=1/x-x+a\ln x\)有两个极值点\(x_1,x_2\)。证明\(\bigl(f(x_1)-f(x_2)\bigr)/(x_1-x_2)<a-2\)。可先利用驻点方程，再选择对数的有理界。

（6）（2021年新高考I卷第22题再研究）回看习题5.28第（7）题。令\(u=1/a,v=1/b\)，不妨设\(u<v\)，再令\(t=v/u\)。将\(u+v\)表示为\(t\)的函数，用本小节的\(\ln t>2(t-1)/(t+1)\)研究其单调性与端点极限，重新证明原结论。
\end{exercise}
